\documentclass[11pt]{article}
\usepackage[margin=1in]{geometry}
\usepackage{xcolor}
\usepackage{booktabs}
\usepackage[hidelinks]{hyperref}
\definecolor{accent}{HTML}{8A1538}
\setlength{\parindent}{0pt}
\setlength{\parskip}{0.5em}

\begin{document}

% ---------- Front page ----------
\begin{titlepage}
\begin{center}

{\large\scshape Northgate University}\\
{\small School of Engineering and Applied Science}\\[0.3em]
{\color{accent}\rule{0.55\textwidth}{1pt}}\\[0.2em]
{\small Department of Mechanical Engineering}

\vspace{2.4cm}

{\small\scshape Capstone Design Project, Final Report}\\[0.8em]
{\Huge\bfseries\color{accent} SolarFlow}\\[0.5em]
{\LARGE A Low-Cost Solar Thermal Water Heater\\[0.2em] for Off-Grid Cabins}

\vspace{2.2cm}

{\large\bfseries\color{accent} Project Team}\\[0.6em]
\begin{tabular}{l l l}
\toprule
\textbf{Name} & \textbf{Student ID} & \textbf{Role} \\
\midrule
Maya Thompson & 20218854 & Team Lead, Thermal Design \\
Daniel Osei & 20221407 & Structures and Mounting \\
Priya Raman & 20219932 & Controls and Instrumentation \\
Lucas Berg & 20224581 & Testing and Data Analysis \\
\bottomrule
\end{tabular}

\vspace{1.8cm}

\begin{tabular}{r l}
\textbf{Supervisor} & Prof. Elena Vasquez \\
\textbf{Second Reader} & Dr. Samuel Kims \\
\textbf{Course} & MECH 4900, Capstone Design \\
\textbf{Submitted} & 30 April 2026 \\
\end{tabular}

\vfill

{\small Submitted in partial fulfilment of the requirements for the degree of\\
Bachelor of Science in Mechanical Engineering}

\end{center}
\end{titlepage}

% ---------- Body ----------
\section{Project Overview}
SolarFlow is a passive thermosiphon water heater designed for off-grid cabins
where grid power is unavailable and propane delivery is costly. The system
uses a \(2\,\mathrm{m}^2\) flat-plate collector, an insulated 120 litre tank,
and no pumps or electronics, keeping the bill of materials under 400 dollars.

\section{Objectives}
The project set three measurable targets: heat 100 litres of water from 15 to
50 degrees Celsius on a clear day, survive a 40 millimetre hail impact test,
and be assembled by two people with hand tools in under four hours.

\section{Outcome Summary}
The final prototype reached 52.4 degrees Celsius in six hours of insolation,
passed the hail impact test with no glazing damage, and was assembled in three
hours and twenty minutes during the demonstration day. All three targets were
met, and the design files have been released under an open hardware licence.

\end{document}
